% 361_Kagome_FFGFT_Vergleich_En_ch.tex
% Johann Pascher, 12 Sept. 2026
% Structural parallels between kinetic kagome frustration
% (Pei et al., PRL 137, 106702, 2026) and FFGFT

\providecommand{\BM}{\textbf{[B]}}
\providecommand{\KM}{\textbf{[K]}}
\providecommand{\SM}{\textbf{[S]}}
\providecommand{\SetzM}{\textbf{[ASSUMPTION]}}
\providecommand{\QM}{\textbf{[Q]}}

\chapter*{The Kagome Lattice and FFGFT: Four Structural Parallels}
\addcontentsline{toc}{chapter}{The Kagome Lattice and FFGFT: Four Structural Parallels}

\begin{abstract}
Pei, Chen, Castelnovo and Moessner (PRL~137, 106702, 31~August 2026, Open
Access) study the singly doped infinite-$U$ Hubbard model on the kagome
lattice. They find resonating-valence-bond (RVB) polarons, disorder-free
self-trapping with band mass increasing by six orders of magnitude, and a
crossover to $\sqrt{3}\times\sqrt{3}$ order in the unpolarised limit, which
they benchmark against the three-state Potts model. This document identifies
four structural parallels to FFGFT and supports each with the wording of the
source. The parallels share one algebraic root: the Frobenius automorphism
$\phi: x\mapsto x^3$ of order~3 on finite fields. A separate section states
what does \emph{not} correspond. Status of all four parallels: \SM\
(structural correspondence established, physical identification open).
\end{abstract}

% ============================================================
\section*{Status Markers}
% ============================================================
\begin{center}
\begin{tabular}{@{}lp{10cm}@{}}
\textbf{[B]} & Algebraically proved. \\[3pt]
\textbf{[K]} & Numerically confirmed. \\[3pt]
\textbf{[S]} & Structurally plausible; physical identification open. \\[3pt]
\textbf{[Q]} & Taken from the source (Pei et al.), quoted verbatim. \\[3pt]
\textbf{[ASSUMPTION]} & Axiomatic postulate. \\
\end{tabular}
\end{center}

\noindent All verbatim quotations are from Pei et al.\ (2026), Open Access
under CC-BY; section references follow the structure of that paper.

% ============================================================
\section{The Source: What Pei et al.\ Show}
% ============================================================

\subsection{Model and question}

The model is the infinite-$U$ Hubbard model on the kagome lattice with a
single hole:
\begin{quote}\QM
``To gain deeper insight into the role of hole kinetics in determining
magnetism in highly frustrated doped Mott insulators, we consider the
single-hole counter-Nagaoka problem on the kagome lattice, using magnetization
as a tuning parameter.'' (Abstract)
\end{quote}
The term \emph{counter-Nagaoka} matters: Nagaoka's theorem predicts
ferromagnetism from the kinetic energy of a hole on bipartite lattices. On
the kagome lattice the kinetics is \emph{frustrated}, and antiferromagnetic
correlations appear instead:
\begin{quote}\QM
``When said kinetic energy is frustrated, it was later shown that
antiferromagnetic correlations can also appear.'' (Introduction)
\end{quote}

\subsection{The physics of a single triangle}

The mechanism is visible on a single triangle. Pei et al.\ describe it:
\begin{quote}\QM
``Ferromagnetic spins have lowest energy $E_0=-t$ due to destructive
interference of the hole motion around the triangle; whereas antiferromagnetic
spins can form a singlet and effectively thread a $\pi$ flux through the
triangle, relieving the frustration and allowing the hole to delocalize and
lower the energy to $E_0=-2t$ (bottom of an unfrustrated 1D tight-binding
band).'' (Section \emph{Model})
\end{quote}
Two points to retain. First: the ferromagnetic configuration is
\emph{destructively interfered} --- the hole cannot move freely. Second: the
singlet \emph{threads a $\pi$ flux} through the triangle. This $\pi$ flux is
not a metaphor but a Berry phase: the hole accumulates phase $\pi$ on one
full circuit, and the interference switches from destructive to
constructive.

\subsection{Flat band and compact localised states}

In the fully polarised case the problem reduces to spinless fermions. The
lowest band is then exactly flat:
\begin{quote}\QM
``The fully polarized case reduces to a well-known spinless fermionic
tight-binding model, where the bottom band is exactly flat (energy $-2$) and
the states are compact localized around individual hexagons.''
(Section \emph{RVB polarons and delocalization})
\end{quote}
Kinetic frustration arises precisely when this flat band is lowest:
\begin{quote}\QM
``Indeed, kinetic frustration arises when the kagome flat band happens to be
lowest in energy for a hole-doped fully polarized electron system close to
half-filling. As such, holes in a fully polarized system occupy compact
localized states at low energy.'' (Introduction)
\end{quote}

\subsection{RVB polarons: formation, delocalisation, self-trapping}

Adding a few reversed spins ($n=1,\ldots,4$) forms polarons that
\emph{delocalise} --- but with strongly enhanced band mass:
\begin{quote}\QM
``Upon adding few reversed spins, polarons form and delocalize, albeit with
a band mass 100 times larger than that of an unfrustrated tight-binding
particle on the same lattice.'' (Introduction)
\end{quote}
The structure of these polarons is described by valence-bond singlets on
tree-like Husimi cacti:
\begin{quote}\QM
``The short distance physics of the kagome lattice is close to tree-like,
and we find that the valence bond language of the same Hamiltonian on the
Husimi cactus proves highly useful to model and understand the polarons in
this Letter.'' (Section \emph{RVB polarons and delocalization})
\end{quote}
At $n=5$ the behaviour changes abruptly:
\begin{quote}\QM
``The polaron bandwidth suddenly drops by at least six orders of magnitude
($<10^{-8}$ in a $4\times4$ system of 48 sites, with 16 near-degenerate
GSs); to the best of our numerics it is indistinguishable from a flat band
of localized states.'' (Section \emph{Strong self-trapping})
\end{quote}
The mechanism of this self-trapping is a resonance of six symmetry-related
clusters around a hexagonal plaquette:
\begin{quote}\QM
``[\ldots] there are six symmetry-related such clusters around any hexagonal
plaquette of the lattice, and the corresponding six Husimi cactus states are
not orthogonal. [\ldots] the GS wave function takes the form of a symmetric
superposition of all six states, spreading across all 12 sites around the
hexagonal plaquette.'' (Section \emph{Strong self-trapping})
\end{quote}
And the interpretation:
\begin{quote}\QM
``By resonating around a hexagonal plaquette, it seems that the 1H5S RVB
polaron acquires a possibly infinite mass, undergoing localization
reminiscent to the compact-localized hexagonal states of the spinless
single-particle problem.'' (Section \emph{Strong self-trapping})
\end{quote}
Notably, this localisation occurs in a \emph{disorder-free} system: the
authors speak of states that ``relocalize in a disorder-free system''
(Conclusions).

\subsection{The unpolarised limit and the Potts model}

As $n$ grows, the valence-bond language loses its usefulness and
semiclassical correlations emerge:
\begin{quote}\QM
``We see that this `frustration on frustration' leads to the appearance of
magnetic correlations consistent with semiclassical $\sqrt{3}\times\sqrt{3}$
order as we approach the unpolarized limit.''
(Section \emph{Unpolarized state})
\end{quote}
The benchmark is the three-state Potts model:
\begin{quote}\QM
``The data presented in Fig.~5 include our Letter of the three-state Potts
model, with Hamiltonian $H_\text{Potts}=\sum_{\langle ij\rangle}
\delta_{\sigma_i\sigma_j}$, $\sigma_i=0,1,2$. The ground state ensemble is
extensively degenerate and consists of all states without identical nearest
neighbors.'' (Appendix \emph{Potts model})
\end{quote}
This model is not chosen arbitrarily --- it marks a phase transition:
\begin{quote}\QM
``This model is useful in that it sits at the Kosterlitz-Thouless (KT)
transition separating a $\sqrt{3}\times\sqrt{3}$ ordered and an
algebraically correlated phase, and it can thus provide a natural benchmark
for an analysis of semiclassical correlations.''
(Section \emph{Unpolarized state})
\end{quote}

% ============================================================
\section{Parallel 1: Frobenius Orbits --- the Same Algebraic Structure}
% ============================================================

\subsection{FFGFT side}

\BM\ (Doc.~339): The multiplicative group $GF(27)^*\cong\mathbb{Z}_{26}$
decomposes under the Frobenius automorphism $\phi: x\mapsto x^3$ into
\begin{itemize}
  \item two fixed points $\{+1,-1\}$ --- the massive $\mathbb{Z}_2$ sector
        (particle/antiparticle),
  \item four three-element orbits in $\mathbb{Z}_{13}$, each
        $\{a,3a,9a\pmod{13}\}$: $\{1,3,9\}$, $\{2,5,6\}$, $\{4,10,12\}$,
        $\{7,8,11\}$,
  \item combined with the two $\mathbb{Z}_2$ parities: eight gluons of the
        adjoint $SU(3)$ representation.
\end{itemize}
The three-element orbits correspond to the three colour states $\{R,G,B\}$;
the parity distinguishes colour from anticolour (Doc.~339, §3.2).

\subsection{Kagome side}

The kagome lattice has three sublattice sites per unit cell; the unit cell
itself is a triangle. Pei et al.\ explicitly use system sizes that are
multiples of~3:
\begin{quote}\QM
``$\sqrt{3}\times\sqrt{3}$ order requires $L_x$ and $L_y$ to be integer
multiples of 3 to avoid incommensurability issues.''
(Appendix \emph{Methods})
\end{quote}
The order itself shows at the K points of the Brillouin zone:
\begin{quote}\QM
``In all cases, we find peaks at the $K$ points, suggestive of
$\sqrt{3}\times\sqrt{3}$ order.'' (Section \emph{Unpolarized state})
\end{quote}
The K points carry threefold rotational symmetry ($C_3$); the structure
factor $S(\mathbf{k})$ is invariant under rotation by $2\pi/3$.

\subsection{Structural correspondence}

Both sides carry the same algebra: a cyclic group on which an automorphism
of order~3 acts and generates three-element orbits. In FFGFT it is $\phi$ on
$GF(27)^*$; on the kagome lattice it is $C_3$ on the sublattice sites and ---
in reciprocal space --- on the K points. \SM

What is \emph{not} claimed: that the four $\mathbb{Z}_{13}$ orbits are
physically identical to four kagome configurations. The kagome lattice has a
single $C_3$ orbit per unit cell; FFGFT has four orbits in $\mathbb{Z}_{13}$.
The correspondence lies in the \emph{order} of the acting automorphism, not
in the number of orbits.

% ============================================================
\section{Parallel 2: Three-State Potts Model = $GF(3)$ = Fixed Field}
% ============================================================

\subsection{FFGFT side}

\BM\ (Doc.~339, §3.3): The fixed field of $\phi$ on $GF(27)$ is
$GF(3)=\{0,+1,-1\}$ (Fermat's little theorem: $x^3=x$ for all $x\in GF(3)$).
In FFGFT, $GF(3)$ corresponds to the $U(1)$ sector: the projection
$GF(27)\to GF(3)$ describes the electrically neutral sector (photon). The
three elements $\{0,+1,-1\}$ are also the three $\mathbb{Z}_3$ sectors of
Doc.~336 (fixed points $\{0,+1,-1\}$ = FFGFT sectors \BM).

\BM\ (Doc.~321, §1.2): The triality condition $T_R=0\pmod{3}$ is equivalent
to confinement:
\begin{quote}
``Confinement is equivalent to requiring all observable states to have
triality $T_R=0$.'' (Doc.~321)
\end{quote}
A single quark ($T_R=1$) is not observable; a meson ($T_R=1+2=0$) and a
baryon ($T_R=1+1+1=0$) are.

\subsection{Kagome side}

The Potts benchmark has exactly three states $\sigma\in\{0,1,2\}$ with
$\mathbb{Z}_3$ symmetry (quoted above). The ground-state condition ``no
identical nearest neighbors'' is a frustration condition: on a triangle all
three states must occur --- exactly one $\mathbb{Z}_3$-neutral configuration
per triangle. The KT transition separates the phase in which this
$\mathbb{Z}_3$ order is long-ranged ($\sqrt{3}\times\sqrt{3}$) from the
algebraically correlated phase.

\subsection{Structural correspondence}

The identification $\sigma\in\{0,1,2\}\leftrightarrow GF(3)=\{0,+1,-1\}$ is
algebraically exact: both are $\mathbb{Z}_3$ as an additive group. The Potts
frustration condition (all three states on every triangle) and the FFGFT
confinement condition (sum of trialities $\equiv0\pmod{3}$) are the same
statement: only $\mathbb{Z}_3$-neutral configurations are admissible. \SM

The Potts KT transition therefore has an FFGFT counterpart: the transition
between confined ($T_R=0$ enforced at long range) and deconfined (triality
only algebraically correlated). Whether this identification holds
physically beyond the algebra is open (§7).

% ============================================================
\section{Parallel 3: Mass as Frozen Kinetic Energy}
% ============================================================

\subsection{FFGFT side}

\BM\ (Doc.~286): The duality $\tilde{T}\cdot m=1$ enforces two energy types:
\begin{center}
\begin{tabular}{lll}
\toprule
 & Static sector & Kinetic sector \\
\midrule
Form        & $E_0=mc^2$ (rest mass)        & $E=hf$ (free) \\
Seat        & $\mathbb{Z}_3$ core (winding) & window \\
Time        & $\tilde{T}=1/m$               & $\tilde{T}=1/\omega$ \\
Status      & derived (Doc.~285)            & open (Doc.~286) \\
\bottomrule
\end{tabular}
\end{center}
Doc.~286 states: ``The rest mass is static: a topologically stable winding
number on $T^4$ [\ldots] \emph{Mass is frozen kinetic energy.}''
\BM\ (Doc.~339): The massive/massless separation is algebraically forced by
the Frobenius --- fixed points $\{+1,-1\}$ are massive, the three-element
orbits are massless (gluons).

\subsection{Kagome side}

The kagome system shows \emph{the same dichotomy}, but not as a sector
separation --- as a \emph{non-monotonic band mass} along the polarisation
axis:
\begin{itemize}
  \item $n=0$: flat band, compact localised --- ``effectively infinite mass''
        (kinetic energy frozen by destructive interference).
  \item $n=1,\ldots,4$: delocalisation with band mass $\times100$
        (``polarons form and delocalize'').
  \item $n=5$: relocalisation, band mass $\times10^6$ (``acquires a possibly
        infinite mass'').
  \item $n\to$ unpolarised: delocalisation (``delocalized, dispersive holes
        are recovered'').
\end{itemize}
Pei et al.\ summarise:
\begin{quote}\QM
``Our Letter uncovers an intricate interplay between single-particle flat
bands and cooperative `polaronic' self-trapping physics.'' (Introduction)
\end{quote}

\subsection{Structural correspondence}

In both systems, \emph{mass} results from geometric freezing of kinetic
freedom: in FFGFT through topological winding on $T^4$ (Doc.~314: ``winding
numbers are topological''), on the kagome lattice through destructive
interference around plaquettes. Both are localisation \emph{without
disorder} --- a purely geometric effect. \SM

The difference: FFGFT has two \emph{sectors} (massive/massless); the kagome
system has one \emph{parameter} (polarisation) along which band mass switches
between the two extremes. The kagome system thus realises the transition
that FFGFT treats as a sector boundary.

% ============================================================
\section{Parallel 4: $\pi$-Flux = Berry Phase = Winding Number}
% ============================================================

\subsection{Kagome side}

The quotation from §1.2 is central here: the singlet ``effectively
thread[s] a $\pi$ flux through the triangle''. The $\pi$ flux is the phase
the hole picks up on one circuit of the triangle. Without it (ferromagnetic)
the hopping amplitude interferes destructively; with it, constructively.
The energy gain is a factor~2 ($E_0=-t\to-2t$).

The same logic repeats at larger scale in the self-trapping: the six Husimi
states around the plaquette interfere \emph{symmetrically} --- constructive
interference for the bound component, spreading the state over all twelve
sites, but destructive for any motion \emph{out} of the plaquette.

\subsection{FFGFT side}

\BM\ (Doc.~314, Ch.~D2): ``in the rolled-up state there are no fractal
corrections --- winding numbers are topological, the defect $d=100\xi$
accumulates only when unrolled.'' Winding numbers on $T^4$ are quantised
Berry phases around closed loops: phase $2\pi n$, $n\in\mathbb{Z}$, on one
circuit of a torus cycle.

\BM\ (Doc.~321, §2.1): The $\mathbb{Z}_3$ generator $\tau$ acts as a
unitary operator on $L^2(T^4)$ with $\tau^3=\mathrm{id}$; eigenvalues are
$\omega^k$, $\omega=e^{2\pi i/3}$. The three sectors $P_k$ are distinguished
by the phase accumulated under a $\mathbb{Z}_3$ rotation.

\BM\ (Doc.~230): Entanglement is ``holonomy around cycles'' --- the phase
accumulated on a closed circuit is what constitutes entanglement
geometrically.

\subsection{Structural correspondence}

The $\pi$ flux in the kagome triangle (Berry phase $\pi$, half winding) and
the $\mathbb{Z}_3$ phases $\omega^k$ in FFGFT (Berry phase $2\pi k/3$) are
the same operation: phase accumulated along a closed path that changes the
accessible Hilbert space. On the kagome lattice it lifts the frustration; in
FFGFT it selects the sector. \SM

% ============================================================
\section{What Does Not Correspond}
% ============================================================

Honesty requires naming the limits of the parallels.

\paragraph{Dimension.} The kagome lattice is two-dimensional; $T^4/Z_3$ is
four-dimensional. The $C_3$ symmetry of the kagome lattice lives in the
plane; the $\mathbb{Z}_3$ action on $T^4$ has nine fixed points (Doc.~314)
and acts on a compact space without boundary. There is no map embedding the
kagome lattice into $T^4/Z_3$.

\paragraph{Classical versus quantum.} The Potts model is classical; Pei et
al.\ use it as a \emph{semiclassical} reference and scale correlators by
$(S^z)^2=1/4$ instead of $S(S+1)=3/4$. FFGFT works on
$L^2(T^4)\otimes\mathbb{C}^3$ (Docs.~230/282) fully quantum-mechanically.
The correspondence in Parallel~2 is algebraic ($\mathbb{Z}_3$), not
dynamical.

\paragraph{Number of orbits.} FFGFT has four three-element orbits in
$\mathbb{Z}_{13}$ (eight gluons); the kagome lattice has one $C_3$ orbit per
unit cell. Parallel~1 concerns the \emph{order} of the automorphism, not the
orbit count.

\paragraph{Structure versus dynamics.} All four parallels are structural.
None predicts that an FFGFT computation reproduces the kagome spectrum or
vice versa. They state: the same algebra ($\phi$ of order~3) generates the
same kind of constraint in both systems (only $\mathbb{Z}_3$-neutral
configurations admissible; localisation without disorder; Berry phase as
selector).

\paragraph{One unexplained number.} The self-trapping cluster spans twelve
sites around a hexagonal plaquette. In Doc.~314 the first excitation of the
$D_4$ lattice has degeneracy $24=12+12$ (``FCC half and winding half
separately addressable''). Whether the~12 of the kagome cluster relates to
one of the 12-halves of $D_4$ is \emph{not} investigated and is not listed
here as a parallel. \SM

% ============================================================
\section{Summary and Open Questions}
% ============================================================

\begin{center}
\begin{tabular}{p{4.2cm}p{4.2cm}p{3.6cm}l}
\toprule
Kagome system (Pei et al.) & FFGFT & Shared algebra & Status \\
\midrule
K-point order, $C_3$ on sublattice &
  Frobenius orbits on $GF(27)^*$ (Doc.~339) &
  automorphism of order~3 & \SM \\
Three-state Potts, KT transition &
  $GF(3)$ fixed field, triality $T_R=0$ (Docs.~339, 321) &
  $\mathbb{Z}_3$ neutrality & \SM \\
Flat band / self-trapping, band mass $\times10^6$ &
  static/kinetic, frozen winding (Docs.~286, 314) &
  localisation without disorder & \SM \\
$\pi$ flux through triangle &
  winding number, $\mathbb{Z}_3$ phase $\omega^k$ (Docs.~314, 321) &
  Berry phase around loop & \SM \\
\bottomrule
\end{tabular}
\end{center}

All four parallels are structurally established \SM. The physical question
--- whether the kagome system is a low-energy realisation of the FFGFT Galois
structure --- remains open and is not claimed here.

\paragraph{Open questions \SM.}
\begin{enumerate}
  \item Is the Potts KT transition (separation of $\mathbb{Z}_3$-ordered /
        algebraically correlated) identifiable with the FFGFT confinement
        transition ($T_R=0$ enforced / merely correlated) at finite
        temperature?
  \item Pei et al.\ announce for two holes ``polaron-mediated interactions
        [\ldots] and possible binding, pairing''. Two holes each with
        $T_R\neq0$ could form a $T_R=0$ state in FFGFT language. Is hole-pair
        binding on the kagome lattice a meson analogue?
  \item The twelve sites of the self-trapping cluster versus $24=12+12$
        from Doc.~314 --- coincidence or structure?
\end{enumerate}

\paragraph{Source.} Pei,~Y.; Chen,~S.~A.; Castelnovo,~C.; Moessner,~R.:
\textit{Kinetic Kagome Magnetism: From Self-Trapping RVB Polarons to
Semiclassical Correlations.}
Phys.~Rev.~Lett.~\textbf{137}, 106702, published 31~August 2026.
DOI: 10.1103/zbxv-vtq8. Open Access (Editors' Suggestion).
